%% LyX 1.2 created this file.  For more info, see http://www.lyx.org/.
%% Do not edit unless you really know what you are doing.
\documentclass[oneside,english]{amsart}
\usepackage[T1]{fontenc}
\usepackage[latin1]{inputenc}

\makeatletter

%%%%%%%%%%%%%%%%%%%%%%%%%%%%%% LyX specific LaTeX commands.
\providecommand{\LyX}{L\kern-.1667em\lower.25em\hbox{Y}\kern-.125emX\@}

%%%%%%%%%%%%%%%%%%%%%%%%%%%%%% Textclass specific LaTeX commands.
 \theoremstyle{plain}    
 \newtheorem{thm}{Theorem}[section]
 \numberwithin{equation}{section} %% Comment out for sequentially-numbered
 \numberwithin{figure}{section} %% Comment out for sequentially-numbered
 \theoremstyle{definition}
  \newtheorem{example}[thm]{Example}

%%%%%%%%%%%%%%%%%%%%%%%%%%%%%% User specified LaTeX commands.
\usepackage{epic,eepic,epsfig}

\usepackage{babel}
\makeatother
\begin{document}

\title{Lectures 13 \& 14}

\maketitle

\section{$M$/$G$/$1$ Queue}


\subsection*{Pollaczek-Khinchin formula}

If service times are general we cannot use birth death processes to
analyse the system. 

Instead we derive the Pollaczek-Khinchin formula. This formula gives
an expression for the expected waiting time in the queue for any distribution
of service times. \[
\bar{W}=\frac{\lambda \left\langle X^{2}\right\rangle }{2(1-\rho )}\]
 where $\bar{W}$ is expected waiting time, $\left\langle X^{2}\right\rangle =\left\langle X^{2}\right\rangle _{p(X)}$,
$\rho =\lambda \bar{X}$ and $\bar{X}=\left\langle X\right\rangle _{p(X)}$. 


\subsection*{Proof of Pollaczek-Khinchin Formula}

We assume that we have a FIFO queue, \emph{i.e.} customers are served
in the order they arrive. We also assume that arrival times are independent
of service times. We will utilise a new concept known as Mean Residual
Service Time. 


\subsubsection*{The Residual Service Time}

Some definitions: 

\begin{tabular}{rcl}
 $W_{i}$&
=&
 Waiting time in queue of $i$th customer.\\
 $R_{i}$&
=&
 `Residual service time' seen by the $i$th customer.\\
 $X_{i}$&
=&
 Service time of the $i$th customer.\\
 $k_{i}$&
=&
 Number of customers in the system when the $i$th customer arrives. \\
\end{tabular}

%
\begin{figure}
\begin{center}\input{diagrams/residualServiceOverview.eepic}\end{center}


\caption{Graphical representation of residual service time.}
\end{figure}


From these definitions we may immediately write down \[
W_{i}=R_{i}+\sum _{j=i-(k_{i}-1)}^{i-1}X_{j}\]
 summing both sides across $i$ we obtain \[
\sum _{i=1}^{I}W_{i}=\sum _{i=1}^{I}R_{i}+\sum _{i=1}^{I}\sum _{j=i-(k_{i}-1)}^{i-1}X_{j}\]
 dividing through by $I$ gives us \[
\sum _{i=1}^{I}\frac{W_{i}}{I}=\sum _{i=1}^{I}\frac{R_{i}}{I}+\sum _{i=1}^{I}\sum _{j=i-(k_{i}-1)}^{i-1}\frac{X_{j}}{I}\]
 in the limit as $I\rightarrow \infty $ this becomes \[
\lim _{I\rightarrow \infty }\sum _{i=1}^{I}\frac{W_{i}}{I}=\lim _{I\rightarrow \infty }\sum _{i=1}^{I}\frac{R_{i}}{I}+\lim _{I\rightarrow \infty }\sum _{i=1}^{I}\frac{(k_{i}-1)X_{i}}{I}\]
assuming time averages are equivalent to steady state averages this
gives \[
\left\langle W\right\rangle =\left\langle R\right\rangle +\left\langle \left(k-1\right)X\right\rangle \]
 Now since the length of the queue as seen by the $i$th arrival is
independent of the $i$th arrival's length this means $\left\langle (k-1)X\right\rangle =\left\langle k-1\right\rangle \left\langle X\right\rangle $
and since $(k-1)$ is the length of the queue we may write \[
\bar{W}=\bar{R}+\bar{N}_{Q}\bar{X}\]
 Using Little's theorem \[
\bar{N}_{Q}=\lambda \bar{W}\]
 we can obtain \[
\bar{W}=\bar{R}+\lambda \bar{X}\bar{W}\]
 and the fact that $\rho =\lambda \bar{X}$ gives us \[
\bar{W}=\bar{R}+\rho \bar{W}\]
 Leading to \begin{equation}
\bar{W}=\frac{\bar{R}}{1-\rho }\label{eqn:arrivals}\end{equation}
 It remains to calculate $\bar{R}$. 


\subsubsection*{Residual Service Time from Departures Perspective}

Plot of the residual service time over time. Here $M\left(t\right)$
is the number of service completions in the interval $[0,\, t]$.

%
\begin{figure}
\begin{center}\input{./diagrams/residualTime.eepic}\end{center}


\caption{Graphical computation of mean residual time.}
\end{figure}


The time average of the mean residual service time can be calculated
in terms of the area under the curve, $A(t)$ which in turn is the
area under each of the triangles which may be written \[
A(t)=\frac{1}{2}\sum _{i=1}^{M(t)}X_{i}^{2}\]
 The average service time is the area divided by time, $t$. \begin{eqnarray}
\frac{A(t)}{t} & = & \frac{M(t)}{t}\frac{A(t)}{M(t)}\\
 & = & \frac{1}{2}\frac{M(t)}{t}\frac{\sum _{i=1}^{M(t)}X_{i}^{2}}{M(t)}
\end{eqnarray}
 As $t\rightarrow \infty $ we find $M(t)\rightarrow \infty $ and
\[
\bar{R}=\frac{1}{2}\lambda \left\langle X^{2}\right\rangle .\]
 If we assume that time averages and ensemble averages are equal we
may substitute back into eqn (\ref{eqn:arrivals}) and obtain the
Pollaczek-Khinchin formula \[
\bar{W}=\frac{\lambda \left\langle X^{2}\right\rangle }{2(1-\rho )}\]


This implies that the total waiting time is \[
\bar{T}=\bar{X}+\frac{\lambda \left\langle X^{2}\right\rangle }{2(1-\rho )}\]
 and through Little's formula we can obtain the expected queue length,
$\bar{N}_{Q}$, and the expected number in the system, $\bar{N}$. 

\begin{example}
Recall that for an $M$/$M$/$1$ system the service times are exponentially
distributed,
\end{example}
\[
P\left(X\right)=\mu \exp \left(-\mu X\right)\]
this means that\[
\bar{X}=\frac{1}{\mu }\]
 then since the variance of the system is $\frac{1}{\mu ^{2}}$\[
\left\langle X\right\rangle ^{2}=\frac{2}{\mu ^{2}}\]
Substituting into the Pollaczek-Khinchin formula we obtain \begin{eqnarray*}
\bar{W} & = & \frac{\rho }{\mu \left(1-\rho \right)}\\
 & = & \frac{\rho }{\mu -\lambda }
\end{eqnarray*}


Deterministic service times --- each service takes the same amount
of time\[
P\left(X\right)=\delta \left(X-\frac{1}{\mu }\right)\]
where $\delta \left(\cdot \right)$is the Dirac delta function. For
this system \[
\left\langle X\right\rangle =\frac{1}{\mu }\]
and \[
\left\langle X^{2}\right\rangle =\frac{1}{\mu ^{2}}\]
\emph{i.e.} the variance is zero as we would expect. Substituting
into the Pollaczek-Khinchin formula we obtain\[
\bar{W}=\frac{\rho }{2\mu \left(1-\rho \right)}\]


In fact we may derive the full range of results for these distributions. 

\begin{center}\begin{tabular}{c|c|c|c}
 Expected value &
 General &
 Exponential &
Deterministic\\
\hline
$\bar{W}$&
 $\frac{\lambda }{2(1-\rho )}\left\langle X^{2}\right\rangle $&
 $\frac{\rho }{\mu (1-\rho )}$&
 $\frac{\rho }{2\mu (1-\rho )}$\\
 $\bar{T}=\bar{W}+\frac{1}{\mu }$&
 $\frac{\lambda }{2(1-\rho )}\left\langle X^{2}\right\rangle +\frac{1}{\mu }$&
 $\frac{1}{\mu (1-\rho )}$&
 $\frac{2-\rho }{2\mu (1-\rho )}$\\
 $\bar{N_{Q}}=\lambda \bar{W}$&
 $\frac{\lambda ^{2}}{2(1-\rho )}\left\langle X^{2}\right\rangle $&
$\frac{\rho ^{2}}{(1-\rho )}$&
 $\frac{\rho ^{2}}{2(1-\rho )}$\\
 $\bar{N}=\lambda \bar{W}+\rho =\lambda \bar{T}$&
 $\frac{\lambda ^{2}}{2(1-\rho )}\left\langle X^{2}\right\rangle +\rho $&
 $\frac{\rho }{1-\rho }$&
 $\frac{\rho (2-\rho )}{2(1-\rho )}$\\
\end{tabular}\end{center}

\begin{example}
Uniform Service times --- A uniform probability distribution is defined
between 0 and $b$ for $X>0$ as\[
U\left(X|b\right)=\left\{ \begin{array}{cc}
 \frac{1}{b} & 0<X\leq b\\
 0 & X>b\end{array}
\right.\]
show that the man of this uniform distribution is given by $\frac{b}{2}$
and its variance by $\frac{b^{2}}{12}$ and find the expected waiting
time for an $M/U/1$ queue?

For the mean we have\begin{eqnarray*}
\left\langle X\right\rangle  & = & \int _{0}^{\infty }Xp\left(X\right)dX\\
 & = & \int _{0}^{\infty }XU\left(X|b\right)dX\\
 & = & \int _{0}^{b}\frac{X}{b}dX\\
 & = & \left[\frac{X^{2}}{2b}\right]_{0}^{b}\\
 & = & \frac{b^{2}}{2b}-0\\
 & = & \frac{b}{2}
\end{eqnarray*}
For the variance we also require \begin{eqnarray*}
\left\langle X^{2}\right\rangle  & = & \int _{0}^{\infty }X^{2}p\left(X\right)dX\\
 & = & \int _{0}^{b}\frac{X^{2}}{b}dX\\
 & = & \left[\frac{X^{3}}{3b}\right]_{0}^{b}\\
 & = & \frac{b^{3}}{3b}-0\\
 & = & \frac{b^{2}}{3}
\end{eqnarray*}
The variance is $\left\langle X^{2}\right\rangle -\left\langle X\right\rangle ^{2}=\frac{b^{2}}{3}-\left(\frac{b}{2}\right)=\frac{b^{2}}{12}$.
\begin{description}
\item [Waiting~Time~For~Queue]For simplicity we consider $\mu =\frac{1}{\left\langle X\right\rangle }.$
So $\mu =\frac{2}{b}$ and thus $b=\frac{2}{\mu }$. So we can write
the second moment in terms of $\mu $. $\left\langle X^{2}\right\rangle =\frac{b^{2}}{3}$,
substituting in $b=\frac{2}{\mu }$ we have $\left\langle X^{2}\right\rangle =\frac{4}{3\mu ^{2}}$.
Substituting in to the P-K formula we recover\[
\bar{W}=\frac{2\rho }{3\mu \left(1-\rho \right)}.\]

\end{description}
\end{example}
~

\begin{example}
Automatic Repeat Request system --- ARQ System. 
\end{example}
%
\begin{figure}
\begin{center}\input{./diagrams/ARQ.eepic}\end{center}


\caption{The ARQ System.}
\end{figure}
Assume each packet takes 1 millisecond to be transmitted.

The system waits $n-1$ milliseconds for an acknowledgement. If there
is no acknowledgement the system retransmits.

The probability of an error in a packet is taken to be $p$. Errors
are taken to be independent. Assume that packets arrive as a Poisson
process with rate $\lambda $.  Time interval between transmission
and re-transmission is $1+kn$ time units with probability $\left(1-p\right)p^{k}$

The probability of the service time is\[
P\left(X=1+kn\right)=\left(1-p\right)p^{k}\]
 The expected service time is therefore\begin{eqnarray*}
\left\langle X\right\rangle  & = & \sum _{k=0}^{\infty }\left(1+kn\right)\left(1-p\right)p^{k}\\
 & = & \left(1-p\right)\left(\sum _{k=0}^{\infty }p^{k}+n\sum _{k=0}^{\infty }kp^{k}\right)\\
 & = & 1+\frac{np}{1-p}
\end{eqnarray*}
 since $\sum _{k=0}^{\infty }p^{k}=\frac{1}{1-p}$ and $\sum _{k=0}^{\infty }kp^{k}=\frac{p}{\left(1-p\right)^{2}}$.
Thus we can write that\[
\rho =\frac{\lambda \left(1-p\right)}{1-p+np}\]
and thus \[
1-\rho =\frac{1-\lambda -p\left(n-1-\lambda \right)}{1-p+np}\]
The second moment of the service time is\begin{eqnarray*}
\left\langle X^{2}\right\rangle  & = & \sum _{k=0}^{\infty }\left(1+kn\right)^{2}\left(1-p\right)p^{k}\\
 & = & \left(1-p\right)\left(\sum _{k=0}^{\infty }p^{k}+2n\sum _{k=0}^{\infty }kp^{k}+n^{2}\sum _{k=0}^{\infty }k^{2}p^{k}\right)\\
 & = & 1+\frac{2np}{1-p}+\frac{n^{2}\left(p+p^{2}\right)}{\left(1-p\right)^{2}}
\end{eqnarray*}
since \[
\sum _{k=0}^{\infty }k^{2}p^{k}=\frac{p+p^{2}}{\left(1-p\right)^{3}}\]
 Substitiuting into the P-K formula gives a very ugly formula!


\section{$M/G/1$ Queues with Vacations }

Consider a server which, when it is free, attends to some other business,
for example it could transmit control and record keeping traffic. 

These periods are sometimes known as vacations --- American for holidays. 

\begin{itemize}
\item If a new arrival arrives during a vacation it must wait until the
vacation is over before being dealt with.
\item If a vacation ends and the server is idle it immediately goes on a
new vacation. 
\end{itemize}
~

%
\begin{figure}
\begin{center}\input{./diagrams/vacationMG1.eepic}\end{center}


\caption{Schematic of $M/G/1$ queue}
\end{figure}


We denote the length of these successive vacations to be $V_{1}$,
$V_{2}$, $\dots $. 

We assume that $V_{1},$$V_{2}$, $\dots $ are i.i.d. random variables
which are also independent of the customer times and the service times. 

For a Poisson arrival system with general service times we will look
to compute the average waiting time for a system of this type. 

%
\begin{figure}
\begin{center}\input{./diagrams/vacationResidual.eepic}\end{center}


\caption{Computation of residual service time for system with vacations.}
\end{figure}


For this system \begin{eqnarray}
\frac{1}{t}\int _{0}^{t}r\left(\tau \right)d\tau  & = & \frac{1}{t}\sum _{i=1}^{M\left(t\right)}\frac{1}{2}X_{i}^{2}+\frac{1}{t}\sum _{i=1}^{L\left(t\right)}\frac{1}{2}V_{i}^{2}\nonumber \\
 & = & \frac{M\left(t\right)}{t}\frac{1}{2}\frac{\sum _{i=1}^{M\left(t\right)}X_{i}^{2}}{M\left(t\right)}+\frac{L\left(t\right)}{t}\frac{1}{2}\frac{\sum _{i=1}^{L\left(t\right)}V_{i}^{2}}{L\left(t\right)}\label{eq:vacationResidual}
\end{eqnarray}


On average the proportion of time available for utilisations is $\left(1-\rho \right)$.
If the average length of a vacation at time $t$ is $\bar{V}_{t}$
then the number of vacations that will have arrived is $L\left(t\right)=\frac{t\left(1-\rho \right)}{\bar{V}_{t}}$.
Therefore if we take the limit of (\ref{eq:vacationResidual}) as
$t\rightarrow \infty $ we find \[
\bar{R}=\frac{1}{2}\lambda \left\langle X^{2}\right\rangle +\frac{\left(1-\rho \right)\left\langle V^{2}\right\rangle }{2\bar{V}}\]


So given that the assumptions that we used to develop the formula 

\[
\bar{W}=\frac{\bar{R}}{1-\rho }\]
are still valid, we may write\[
\bar{W}=\frac{\lambda \left\langle X^{2}\right\rangle }{2\left(1-\rho \right)}+\frac{\left\langle V^{2}\right\rangle }{2\bar{V}}\]

\end{document}
